\section{Описание}

\subsection{Источники данных}
Для построения корпуса использовались два источника данных.

\textbf{Подкорпус текстов песен.} Список популярных треков получается из
\MYhref{https://www.last.fm/api}{Last.fm API} (метод \texttt{chart.gettoptracks}),
а текст песни для каждого трека --- из бесплатного
\MYhref{https://api.lyrics.ovh}{Lyrics.ovh API} (запрос вида
\texttt{/v1/\{артист\}/\{название\}}). Last.fm используется с бесплатным ключом
API, Lyrics.ovh ключа не требует.

\textbf{Подкорпус информации о записях.} Список популярных артистов получается
из Last.fm API (метод \texttt{chart.gettopartists}), затем для каждого артиста из
\MYhref{https://musicbrainz.org/doc/MusicBrainz_API}{MusicBrainz API} загружается
список записей (поиск \texttt{/ws/2/recording}) и подробная информация о каждой
записи: длительность, жанры, идентификатор MBID.

\subsection{Характеристики корпуса}
Корпус состоит из двух подкорпусов. Каждый документ --- HTML-страница,
сгенерированная программой, плюс отдельный JSON-файл с метаданными.

В подкорпусе текстов песен документ содержит название трека, исполнителя и текст
песни (блок \texttt{div.lyrics}). Метаданные: артист, название, длина текста,
число слов, время загрузки, источник данных.

В подкорпусе записей MusicBrainz документ содержит название записи, исполнителя,
длительность, жанры и MBID. Метаданные: MBID, артист, название, время загрузки.

Разметка документов --- HTML, метаданные хранятся в формате JSON.

\subsection{Разбиение на документы}
Корпус разбит на документы по принципу \enquote{один объект --- один документ}:
каждый трек сохраняется в файл \texttt{track\_\{N\}.html} с метаданными
\texttt{track\_\{N\}\_full.json}, каждая запись MusicBrainz --- в файл
\texttt{recording\_\{mbid\}.html} с метаданными \texttt{recording\_\{N\}.json}.
Такое разбиение позволяет обрабатывать документы независимо друг от друга.

\subsection{Выделение текста}
Сырые данные --- HTML-страницы. Текст выделяется из размеченных блоков: для
текстов песен --- содержимое блока \texttt{div.lyrics}, для записей --- текстовое
содержимое информационных блоков (название, исполнитель, длительность, жанры).
HTML-теги, CSS-оформление и служебные элементы (стили, футер, статистика)
отбрасываются.

\subsection{Существующие поисковики}
Для поиска по выбранному набору документов уже можно использовать:

\begin{itemize}
    \item \textbf{Google с ограничением по сайтам} --- для поиска текстов песен:
    тексты тех же песен доступны на публичных сайтах (genius.com, azlyrics.com),
    поэтому запросы вида \texttt{site:genius.com \{артист\} \{название\}} находят
    те же документы, что и в корпусе;
    \item \textbf{встроенный поиск MusicBrainz}
    (\MYhref{https://search.musicbrainz.org}{search.musicbrainz.org}) --- для
    поиска записей по артисту и названию.
\end{itemize}

\subsection{Примеры запросов и недостатки выдачи}

\begin{enumerate}
    \item Google: \texttt{Olivia Rodrigo "stupid song" site:genius.com} --- текст
    находится, но его версия может отличаться от скачанной с Lyrics.ovh (другая
    пунктуация и форматирование).
    \item Google: \texttt{Taylor Swift "Patient Zero" текст песни} --- в выдаче
    много дублей одного и того же текста на разных сайтах, страницы с рекламой и
    новости, а не только сам текст.
    \item Поиск MusicBrainz: \texttt{artist:"Radiohead" recording:"Just"} ---
    нужная запись находится, но вместе с ней десятки похожих (ремастеры,
    live-версии, сборники), из которых сложно выбрать нужную.
\end{enumerate}

Общие недостатки: поисковики не ограничены нашим корпусом (выдают документы,
которых в нём нет), а поиск MusicBrainz работает только по метаданным, а не по
содержимому документа.

\subsection{Статистика по корпусу}
Статистика приведена для целевого размера корпуса --- 30000 документов (по 15000
в каждом подкорпусе); значения экстраполированы по фактически скачанным
документам.

\begin{tabular}{|l|c|c|c|}
\hline
Показатель & Тексты песен & Записи & Всего \\ \hline
Размер \enquote{сырых} данных & $\sim$57 МБ & $\sim$41 МБ & $\sim$98 МБ \\ \hline
Количество документов & 15000 & 15000 & 30000 \\ \hline
Размер выделенного текста & $\sim$20 МБ & $\sim$6 МБ & $\sim$26 МБ \\ \hline
Средний размер документа & $\sim$3{,}8 КБ & $\sim$2{,}7 КБ & $\sim$3{,}3 КБ \\ \hline
Средний объём текста в документе & $\sim$1{,}3 КБ & $\sim$0{,}4 КБ & $\sim$0{,}9 КБ \\ \hline
\end{tabular}

\pagebreak

\section{Исходный код}
Программа \texttt{parser.py} написана на Python 3 и использует библиотеку
\texttt{requests}. Она состоит из двух классов: \texttt{LyricsOvhParser}
(подкорпус текстов песен) и \texttt{MusicBrainzParser} (подкорпус записей).

Этапы работы \texttt{LyricsOvhParser}:
\begin{enumerate}
    \item получение страницы популярных треков из Last.fm API;
    \item для каждого трека --- запрос текста из Lyrics.ovh;
    \item сохранение HTML-документа и JSON-метаданных.
\end{enumerate}

Этапы работы \texttt{MusicBrainzParser}:
\begin{enumerate}
    \item получение страницы популярных артистов из Last.fm API;
    \item поиск записей артиста в MusicBrainz (постранично, по 100 записей);
    \item загрузка подробной информации о каждой записи;
    \item сохранение HTML-документа и JSON-метаданных.
\end{enumerate}

При написании кода пришлось учесть особенности реальных API: временные сбои
серверов (повторные попытки запроса), зависание отдельных запросов (таймауты),
ограничение частоты запросов MusicBrainz (пауза 1{,}2~с между запросами) и
отсутствие части текстов (около 15--20\% треков возвращают 404, такие треки
пропускаются).

Получение списка популярных треков из Last.fm (с повторными попытками при сбое):

\begin{lstlisting}[language=Python]
def get_popular_tracks(self, page=1, page_size=20):
    api_key = "..."
    url = "https://ws.audioscrobbler.com/2.0/"
    params = {
        "method": "chart.gettoptracks",
        "api_key": api_key,
        "format": "json",
        "page": page,
        "limit": page_size
    }
    for attempt in range(3):
        try:
            response = self.session.get(url, params=params, timeout=30)
            if response.status_code == 200:
                data = response.json()
                tracks = data.get("tracks", {}).get("track", [])
                track_list = []
                for track in tracks:
                    artist = track.get("artist", {}).get("name", "")
                    title = track.get("name", "")
                    if artist and title:
                        track_list.append((artist, title))
                return track_list
        except Exception:
            pass
        time.sleep(2 * (attempt + 1))
    return None
\end{lstlisting}

Запрос текста песни из Lyrics.ovh (таймаут и повтор при ошибке сервера):

\begin{lstlisting}[language=Python]
def get_lyrics(self, artist, title):
    url = f"{self.base_url}/{quote_plus(artist)}/{quote_plus(title)}"
    for attempt in range(2):
        try:
            response = self.session.get(url, timeout=30)
            if response.status_code == 200:
                lyrics = response.json().get("lyrics", "")
                if lyrics and len(lyrics.strip()) > 50:
                    return lyrics.strip()
                return None
            elif response.status_code == 404:
                return None
            elif response.status_code >= 500 and attempt == 0:
                continue
            return None
        except Exception:
            if attempt == 0:
                continue
            return None
    return None
\end{lstlisting}

Основной цикл загрузки текстов (различаются ошибка API и конец списка,
защита от вечной пагинации при отказе сервиса):

\begin{lstlisting}[language=Python]
while downloaded < max_tracks:
    tracks = self.get_popular_tracks(page, page_size)
    if tracks is None:            # API error -> retry the page
        page_errors += 1
        if page_errors >= 3:
            break
        continue
    if not tracks:                # track list is exhausted
        break
    for artist, title in tracks:
        lyrics = self.get_lyrics(artist, title)
        if lyrics:
            self.save_lyrics_html(artist, title, lyrics, downloaded + 1)
            downloaded += 1
        else:
            consecutive_failures += 1
            if consecutive_failures >= 50:
                break
    page += 1
\end{lstlisting}

Основной цикл загрузки записей MusicBrainz (топ-артисты Last.fm, постраничный
поиск записей, пропуск повторов по MBID):

\begin{lstlisting}[language=Python]
while downloaded < max_recordings:
    artists = self.get_top_artists(artist_page, 50)
    if not artists:
        break
    for artist in artists:
        offset = 0
        while downloaded < max_recordings:
            recordings = self.search_recordings_by_artist_api(
                artist, limit=100, offset=offset)
            if recordings is None:    # search failed -> skip artist
                break
            if not recordings:        # no more recordings
                break
            for recording in recordings:
                if recording["id"] in processed_ids:
                    continue
                processed_ids.add(recording["id"])
                data = self.get_recording_details(recording["id"])
                if data:
                    self.save_recording_html(data, recording["id"])
                    downloaded += 1
                time.sleep(1.2)       # MusicBrainz rate limit: 1 req/s
            if len(recordings) < 100:
                break
            offset += 100
    artist_page += 1
\end{lstlisting}

\pagebreak
